\section{Geometric applications}

The following families illustrate the quotient using classical incidence
results. The first has zero defect; the second gives a homocyclic family.

\subsection{Finite M\"obius planes}

A finite M\"obius plane is a
$3$-$(q^{2}+1,q+1,1)$ design; its star differences are $q^{2}-1$ and $q$, with
$\gcd=1$, so $d=1$ and Corollary~\ref{cor:d1} gives the full incidence Smith
form $1^{\,q^{2}}(q+1)$ and $\Delt=0$. This recovers the Miquelian and arbitrary even-order Smith forms
of Ringe~\cite{Ringe2008} from the star-difference bound.

\subsection{Hermitian unitals}

The classical Hermitian unital $H(q)$---the
$2$-$(q^{3}+1,q+1,1)$ design of Key's explicit finite-field/secant-line
construction~\cite{Key}---has, by Hiss's incidence divisors~\cite{Hiss},
the incidence cokernel
\[
  \coker(A(H(q)))\cong(\ZZ/(q+1))^{q^{2}-q+1}.
\]
Removing the augmentation factor via Theorem~\ref{thm:hs} yields the
homocyclic family
\[
  \Delt(H(q))\cong(\ZZ/(q+1))^{q^{2}-q},
\]
so $\Delt(H(4))=(\ZZ/5)^{12}$. The incidence Smith form is Hiss's~\cite{Hiss}; the displayed defect
is its quotient by the augmentation factor.

A nonzero $t=3$ instance---the design $\SQS(8)=\AG(3,2)$, for which $d=2$ and
$\Delt\cong(\ZZ/2)^{3}$---is worked out in Appendix~\ref{sec:app-sqs}, where it
shows that the hypothesis $d=1$ of Corollary~\ref{cor:d1} has content.
